\chapter{Express Routing}

\section{HTTP Methods and Their Purpose}

\begin{itemize}
  \item \textbf{GET} --- retrieve; no body, safe and idempotent.
  \item \textbf{POST} --- create a new resource.
  \item \textbf{PUT} --- replace a resource entirely.
  \item \textbf{PATCH} --- partially update a resource.
  \item \textbf{DELETE} --- remove a resource.
\end{itemize}

\begin{notebox}[NOTE]
Many REST APIs use \texttt{PUT} where \texttt{PATCH} is more correct. This handbook uses \texttt{PUT} for task updates since the form sends the full object.
\end{notebox}

\section{Route Files with express.Router()}

\whatwhyhow{
\texttt{express.Router()} creates a mini, mountable Express app scoped to one resource.
}{
Without it, every route lives in one giant \texttt{app.js}, unmanageable as the API grows.
}{
Define routes in \texttt{routes/*.routes.js}, export the router, mount it in \texttt{app.js} with \texttt{app.use("/api/tasks", taskRoutes)}.
}

\subsection{routes/auth.routes.js}

\begin{lstlisting}[language=JavaScript]
const express = require("express");
const { register, login, logout, getMe } = require("../controllers/auth.controller");
const { protect } = require("../middleware/auth.middleware");
const router = express.Router();

router.post("/register", register);
router.post("/login", login);
router.post("/logout", logout);
router.get("/me", protect, getMe);
module.exports = router;
\end{lstlisting}

\subsection{routes/task.routes.js}

\begin{lstlisting}[language=JavaScript]
const express = require("express");
const { createTask, getTasks, getTaskById, updateTask, deleteTask } = require("../controllers/task.controller");
const { protect } = require("../middleware/auth.middleware");
const router = express.Router();

router.use(protect); // every task route requires authentication
router.post("/", createTask);
router.get("/", getTasks);
router.get("/:id", getTaskById);
router.put("/:id", updateTask);
router.delete("/:id", deleteTask);
module.exports = router;
\end{lstlisting}

\begin{notebox}[NOTE]
\texttt{user.routes.js} and \texttt{profile.routes.js} follow the same router + protect pattern shown above, just scoped to their own controllers.
\end{notebox}

\section{Route Params vs Query Params vs Body}

\begin{itemize}
  \item \textbf{Route params} --- identify a specific resource: \texttt{/tasks/:id} $\rightarrow$ \texttt{req.params.id}.
  \item \textbf{Query params} --- filtering/pagination/sorting: \texttt{/tasks?page=1\&limit=10} $\rightarrow$ \texttt{req.query}.
  \item \textbf{Request body} --- JSON payload for create/update, populated by \texttt{express.json()}.
\end{itemize}

\begin{lstlisting}[language=JavaScript]
// GET /api/tasks?page=2&limit=10&status=todo -> req.query
const { page = 1, limit = 10, status } = req.query;
// GET /api/tasks/:id -> req.params
const { id } = req.params;
\end{lstlisting}

\section{Task Routes Reference}

\begin{table}[h]
\centering
\begin{tabularx}{\textwidth}{l l X}
\toprule
\textbf{Method} & \textbf{Path} & \textbf{Purpose} \\
\midrule
POST   & /api/tasks      & Create task \\
GET    & /api/tasks      & List tasks (paginated/filtered) \\
GET    & /api/tasks/:id  & Fetch one task \\
PUT    & /api/tasks/:id  & Update task \\
DELETE & /api/tasks/:id  & Remove task \\
\bottomrule
\end{tabularx}
\end{table}

\begin{tipbox}[TIP]
Keep route files thin: wire paths to controllers and middleware only, no business logic.
\end{tipbox}
